% !TEX root = ../../main.tex
%% Défi du chapitre 1 — Points fixes et théorème de Cantor-Bernstein
%% Résultats vérifiés par verifications/defis_01_05.py.

\rubrique{Défi}

\begin{defi}{Points fixes et théorème de Cantor-Bernstein}{1 h 30}
Soit $E$ un ensemble. Une application $\varphi:\parties{E}\to\parties{E}$
est dite \emph{croissante} si, pour toutes parties $X,Y$ de $E$,
$X\subseteq Y\implies\varphi(X)\subseteq\varphi(Y)$. Le but du problème est
de montrer qu'une telle application possède toujours un point fixe, puis
d'en tirer un théorème célèbre : s'il existe une injection de $E$ dans $F$
et une injection de $F$ dans $E$, alors $E$ et $F$ sont en bijection.

\partieExo{A}{Le théorème du point fixe de Knaster-Tarski}
Soit $\varphi:\parties{E}\to\parties{E}$ croissante. On pose
\[
\mathcal{A}=\ens{X\in\parties{E}\mid X\subseteq\varphi(X)}
\qquad\text{et}\qquad
A=\bigcup_{X\in\mathcal{A}}X .
\]
\begin{questions}
  \item Montrer que $\mathcal{A}$ n'est pas vide.
  \item Soit $X\in\mathcal{A}$. Montrer que $X\subseteq\varphi(A)$ ; en
    déduire que $A\subseteq\varphi(A)$.
  \item Montrer que $\varphi(A)\in\mathcal{A}$, puis que $\varphi(A)=A$.
  \item Montrer que $A$ est le \emph{plus grand} point fixe de $\varphi$ :
    tout $Y$ vérifiant $\varphi(Y)=Y$ est inclus dans $A$.
  \item L'hypothèse de croissance est-elle utile ? Étudier
    $X\mapsto E\setminus X$ pour $E$ non vide.
\end{questions}

\partieExo{B}{Le théorème de Cantor-Bernstein}
Soient $E$, $F$ deux ensembles, $f:E\to F$ et $g:F\to E$ deux applications
\emph{injectives}. Pour $X\subseteq E$, on pose
$\varphi(X)=E\setminus g\bigl(F\setminus f(X)\bigr)$.
\begin{questions}[resume]
  \item Montrer que $\varphi$ est croissante.
  \item Soit $A$ un point fixe de $\varphi$. Montrer que
    $E\setminus A=g\bigl(F\setminus f(A)\bigr)$.
  \item On définit $h:E\to F$ par $h(x)=f(x)$ si $x\in A$, et
    $h(x)=g^{-1}(x)$ si $x\in E\setminus A$. Justifier que $h$ est bien
    définie, puis montrer que $h$ est bijective.
\end{questions}

\medskip
\begin{center}
\begin{tikzpicture}[line cap=round]
  \foreach \xx/\hh/\bb/\nom in {0/A/E\setminus A/E, 5/f(A)/F\setminus f(A)/F}{%
    \begin{scope}[xshift=\xx cm]
      \begin{scope}
        \clip[rounded corners=3mm] (0,0) rectangle (2.5,3.2);
        \fill[otGreen!16] (0,1.55) rectangle (2.5,3.2);
      \end{scope}
      \draw[otNavy,line width=0.8pt,rounded corners=3mm] (0,0) rectangle (2.5,3.2);
      \draw[otNavy,line width=0.5pt,dash pattern=on 2pt off 1.5pt] (0,1.55) -- (2.5,1.55);
      \node[font=\small,text=otGreenDark] at (1.25,2.38) {$\hh$};
      \node[font=\small,text=otNavy] at (1.25,0.78) {$\bb$};
      \node[font=\small,text=otNavy,above] at (1.25,3.22) {$\nom$};
    \end{scope}}
  \draw[otGreenDark,line width=1pt,-{Stealth[length=2.2mm]}]
    (2.6,2.38) to[bend left=14] node[above,font=\small] {$f$} (4.9,2.38);
  \draw[otNavy,line width=1pt,-{Stealth[length=2.2mm]}]
    (4.9,0.78) to[bend left=14] node[below,font=\small] {$g$} (2.6,0.78);
\end{tikzpicture}\\[0.5mm]
\legende{Le découpage de Cantor-Bernstein : $f$ envoie $A$ sur $f(A)$, et
$g$ envoie le reste de $F$ exactement sur le reste de $E$.}
\end{center}

\partieExo{C}{Applications}
\begin{questions}[resume]
  \item On prend $E=\intFF{0}{1}$, $F=\intFO{0}{1}$, $f(x)=x/2$ et
    $g(y)=y$. Déterminer les points fixes de l'application $\varphi$ de la
    partie B, puis expliciter la bijection $h$ obtenue à partir du plus
    grand d'entre eux. Faire un dessin.
  \item En utilisant les injections $n\mapsto(n,0)$ et
    $(p,q)\mapsto2^{p}3^{q}$ (unicité de la décomposition en facteurs
    premiers, chapitre 7), montrer que $\N$ et $\N\times\N$ sont en
    bijection.
  \item Montrer qu'il n'existe aucune injection de $\parties{E}$ dans
    $E$. \emph{On pourra utiliser le théorème de Cantor (exercice 24).}
\end{questions}
\end{defi}
